\documentclass[aspectratio=169,11pt]{beamer}

\usetheme{Madrid}
\usecolortheme{default}
\usefonttheme{professionalfonts}
\setbeamertemplate{navigation symbols}{}
\setbeamertemplate{footline}[frame number]

\usepackage{amsmath, amssymb, booktabs, array}

\title{Comparative And Global Gender In Labor-Market Development}
\subtitle{Gender Study --- Week 8}
\author{Teaching deck draft}
\date{}

\begin{document}

\begin{frame}
  \titlepage
\end{frame}

\begin{frame}{Central question}
\begin{itemize}
    \item Which gender-and-labor mechanisms travel across countries?
    \item Which depend on development, institutions, norms, and labor-market structure?
    \item How do we use country-specific evidence without overclaiming?
    \item What should readers learn beyond the setting?
\end{itemize}
\end{frame}

\begin{frame}{Why comparative gender research matters}
\begin{itemize}
    \item Gender gaps are labor-market outcomes, not only social indicators.
    \item The same observed gap can reflect supply, demand, care, sectoral structure, formality, mobility, law, or norms.
    \item Cross-country variation helps separate portable mechanisms from institutional transmission.
    \item The practical goal is to make country-specific knowledge legible to labor economists.
\end{itemize}
\end{frame}

\begin{frame}{Comparative taxonomy}
\scriptsize
\begin{tabular}{p{0.25\linewidth} p{0.31\linewidth} p{0.27\linewidth}}
\toprule
Regime or setting & Core mechanism & Gendered labor margin \\
\midrule
High-income welfare states & care policy, tax design, wage setting & hours, retention, promotion \\
Structural transformation & services, informality, migration & participation, sector, formality \\
Low-income/high informality & household production, mobility, legal barriers & entry, self-employment \\
High-migration settings & movement rules, networks, safety & search radius, occupation \\
Restrictive legal/norm settings & family law, social norms, public employment & acceptable jobs, sector sorting \\
Public-employment-heavy settings & state demand, benefits, wage compression & entry, retention, authority \\
\bottomrule
\end{tabular}
\end{frame}

\begin{frame}{Portable mechanism vs setting-specific detail}
\begin{tabular}{p{0.32\linewidth} p{0.48\linewidth}}
\toprule
Object & Interpretation \\
\midrule
Portable mechanism & causal channel that can recur across settings \\
Setting-specific detail & rule, institution, norm, or market structure that shapes transmission \\
Transportability condition & context variable needed for the mechanism to operate elsewhere \\
Policy transferability & whether a specific policy package can be moved across settings \\
\bottomrule
\end{tabular}
\vspace{0.4cm}
\begin{itemize}
    \item A mechanism can travel even when the policy effect size does not.
    \item A policy label can look similar while financing, coverage, and enforcement differ.
\end{itemize}
\end{frame}

\begin{frame}{What it means for a mechanism to travel}
\begin{itemize}
    \item The causal channel is recognizable in another setting.
    \item Relevant context variables are present or can be measured.
    \item The paper states where the mechanism should and should not operate.
    \item The paper does not require the country to ``represent the world.''
\end{itemize}
\vspace{0.3cm}
\[
\text{Design in setting } S \Rightarrow \text{update about mechanism } M
\text{ under conditions } C.
\]
\end{frame}

\begin{frame}{Labor supply, care, and service-sector development}
\begin{itemize}
    \item Labor supply: household production, fertility timing, safety, migration, mobility, norms.
    \item Care infrastructure: childcare, eldercare, school schedules, leave, household help markets.
    \item Service-sector growth: changes labor demand, task mix, hours, and returns to market work.
    \item Comparative question: does demand growth translate into employment when care or mobility constraints bind?
\end{itemize}
\end{frame}

\begin{frame}{Family policy evidence}
\begin{itemize}
    \item Olivetti and Petrongolo: family policies and gender gaps across high-income countries.
    \item Relevant margins: participation, hours, retention, job continuity, wages, and child penalties.
    \item Policy labels are too coarse: coverage, financing, job protection, father eligibility, and provider supply matter.
    \item Mechanism: care constraints and career continuity.
\end{itemize}
\end{frame}

\begin{frame}{Formality, migration, law, and public employment}
\begin{itemize}
    \item Formality changes who is covered by law, social insurance, enforcement, and observed wage data.
    \item Migration and commuting rules shape search radius and feasible labor-market geography.
    \item Legal regimes define formal rights, but enforcement determines credible rights.
    \item Public employment can raise female labor demand while shaping sector sorting and authority.
\end{itemize}
\end{frame}

\begin{frame}{Global evidence map}
\scriptsize
\begin{tabular}{p{0.27\linewidth} p{0.20\linewidth} p{0.22\linewidth} p{0.19\linewidth}}
\toprule
Literature & Setting & Portable mechanism & Detail to respect \\
\midrule
Family policy & high-income & care and continuity & welfare-state package \\
Long-run development & historical & sectoral transformation & timing and composition \\
Social norms & developing countries & acceptable work & norm content \\
Gendered laws & global legal data & formal feasibility & enforcement \\
Informality & development & contract form & informal institutions \\
Global distortions & cross-country & supply/demand wedges & data quality \\
\bottomrule
\end{tabular}
\end{frame}

\begin{frame}{Recent frontier work}
\begin{itemize}
    \item Global distortions: decompose gender inequality into supply-side and demand-side wedges.
    \item External validity: ask how study sites differ from target settings.
    \item Policy portability: state which institutional conditions make a policy effect transportable.
    \item Measurement frontier: combine legal, household, firm, time-use, mobility, and sector data.
\end{itemize}
\end{frame}

\begin{frame}{How to sell country-specific research}
\begin{enumerate}
    \item Identify the labor mechanism.
    \item Map the setting into a comparative taxonomy.
    \item Explain why the setting gives leverage.
    \item Distinguish internal validity, transportability, and external validity.
    \item State what should and should not generalize.
    \item Return to employment, hours, wages, formality, mobility, care, authority, or welfare.
\end{enumerate}
\end{frame}

\begin{frame}{Country-specific evidence is not a liability}
\begin{itemize}
    \item Local institutional knowledge can reveal the real treatment and comparison group.
    \item Language, field access, and administrative detail can make mechanisms observable.
    \item The setting becomes a contribution when it gives leverage on a broader labor object.
    \item The introduction must answer: why should a labor economist who does not know this country care?
\end{itemize}
\end{frame}

\begin{frame}{Common critiques}
\begin{itemize}
    \item Overclaiming generalizability.
    \item Weak mechanism articulation.
    \item Country-as-exotic-case framing.
    \item Weak comparison-group logic.
    \item Inability to connect to broader labor-market objects.
    \item Institutional detail that does not affect identification or transmission.
\end{itemize}
\end{frame}

\begin{frame}{What makes comparative work persuasive}
\begin{itemize}
    \item Starts with a labor object.
    \item Names the mechanism before the country narrative.
    \item Uses comparison groups that map to the mechanism.
    \item Treats institutions as transmission, not decoration.
    \item Tells readers exactly what to update about another setting.
\end{itemize}
\end{frame}

\begin{frame}{Week 8 lab}
\begin{itemize}
    \item Reproduce: synthetic comparative family-policy summaries inspired by Olivetti and Petrongolo.
    \item Diagnose: classify supply, demand, care, legal, and transportability margins.
    \item Transfer: norms, mobility, care, and service-sector demand exercise inspired by Jayachandran.
    \item Optional memo: supply-side and demand-side distortions inspired by Goldberg et al.
\end{itemize}
\end{frame}

\begin{frame}{Research ladder}
\begin{enumerate}
    \item Labor object.
    \item Mechanism.
    \item Setting taxonomy.
    \item Variation.
    \item Comparison.
    \item Transportability conditions.
    \item Contribution.
\end{enumerate}
\end{frame}

\begin{frame}{Bridge to Week 9 and Week 10}
\begin{itemize}
    \item Week 9 asks which data and methods make comparative mechanisms visible.
    \item Week 10 turns the sequence into research designs.
    \item Comparative evidence is strongest when it moves from country knowledge to mechanism, design, and transportability.
\end{itemize}
\end{frame}

\begin{frame}{Takeaways}
\begin{itemize}
    \item Comparative gender research is labor economics when it studies labor objects and mechanisms.
    \item Mechanisms may travel even when institutions and policy effects do not.
    \item Policy transferability requires more than external validity.
    \item Country-specific evidence should be sold through mechanism, taxonomy, leverage, and limits.
\end{itemize}
\end{frame}

\end{document}
